\section{Non-acyclic rank 3 cluster algebras}\label{section4}
In this section, we consider non-acyclic skew-symmetric rank 3 cluster algebras of geometric type.

\subsection[Def\/inition and properties of $\tau$]{Def\/inition and properties of $\boldsymbol{\tau}$}\label{section4.1}

Let  $m\geq 3$ be an integer, $\tilde{B}=(b_{ij})$ be an $m \times 3$ matrix whose principal part $B$ is skew-symmetric and non-acyclic (i.e., $b_{12}$, $b_{23}$, $b_{31}$ are of the same sign).

Def\/ine the map
\begin{gather*}
\tau(B):=(|b_{23}|,|b_{31}|,|b_{12}|)\in \mathbb{Z}^3_{\geq 0}.
\end{gather*}

\begin{Definition}
 Let $(x,y,z) \in \mathbb{R}^3$. We def\/ine a partial ordering ``$\le$'' on  $\mathbb{R}^3$ by $(x,y,z) \leq (x',y',z')$ if and only if $x \leq x'$, $y \leq y'$, and $z \leq z'$.
\end{Definition}

Def\/ine three involutary functions $\mu_1,\mu_2,\mu_3\colon \mathbb{R}^3\to\mathbb{R}^3$ as follows
\begin{gather*}
\mu_1\colon \ (x,y,z)\mapsto (x,xz-y,z), \qquad \mu_2\colon \ (x,y,z) \mapsto  (x,y,xy-z),
\\
\mu_3\colon \  (x,y,z) \mapsto(yz-x,y,z).
\end{gather*}
Let $\Gamma$ be the group generated by $\mu_1$, $\mu_2$, and $\mu_3$.  It follows that the $\Gamma$-orbit of $\tau(B)$ is identical to the set $\{\tau(B')$: $B'$ is in the mutation class of~$B \}$.

 Consider the following two situations.
\begin{enumerate}\itemsep=0pt
\item[(M1)] $(a,b,c)\leq \mu_i(a,b,c)$ for all $i=1,2,3$.
\item[(M2)] $(a,b,c)\leq \mu_i(a,b,c)$ for precisely two indices~$i$.
\end{enumerate}

The following statement follows from \cite[Theorem~1.2, Lemma~2.1]{BBH} for a coef\/f\/icient free cluster algebra. The result holds more generally for a cluster algebra of geometric type since the coef\/f\/icients do not af\/fect whether or not the seed is acyclic.
\begin{Theorem}
\label{cyclic}
Let $B$ be as above, $(a,b,c)=\tau(B)$. Then the following are equivalent:
\begin{itemize}\itemsep=0pt
\item[\rm(i)]$\mathcal{A}({\bf x},{\bf y},B)$ is non-acyclic.

\item[\rm(ii)] $a,b,c \geq 2$ and $abc+4\ge a^2+b^2+c^2$.

\item[\rm(iii)] $a,b,c\ge 2$ and there exists a unique triple in the $\Gamma$-orbit of $\tau(B)$ that satisfies~{\rm(M1)}.
\end{itemize}
Moreover, under these conditions, triples in the $\Gamma$-orbit of $\tau(B)$ not satisfying~{\rm(M1)} must sa\-tis\-fy~{\rm(M2)}.
\end{Theorem}

Assume $\mathcal{A}({\bf x},{\bf y},B)$ is non-acyclic. We call the unique triple in Theorem \ref{cyclic}(iii) the \emph{root} of the $\Gamma$-orbit of~$\tau(B)$.


\begin{Lemma}\label{lemma:chain inequality}
Let $B$ be as above, $(a,b,c)=\tau(B)$.
\begin{enumerate}\itemsep=0pt
\item[$(i)$] The root $(a,b,c)$ of the $\Gamma$-orbit of $\tau(B)$ is the minimum of the $\Gamma$-orbit, i.e., $(a,b,c)\le(a',b',c')$ for any $(a',b',c')$ in the $\Gamma$-orbit.

\item[$(ii)$] If $\mu_k(a,b,c)=(a,b,c)$ for any $k$, then $a=b=c=2$ and $(2,2,2)$ is the unique triple in the $\Gamma-$orbit of~$\tau(B)$. Therefore $(2,2,2)$ must also be the root.

\item[$(iii)$] Assume that $\tau(B)$ is the root of the $\Gamma$-orbit of $\tau(B)$,  $i_1,\dots,i_l\in\{1,2,3\}$, $i_s\neq i_{s+1}$ for $1\le s\le l-1$.
Let $B_0=B$, $B_j=\mu_{i_j}(B_{j-1})$ for $j=1,\dots,l$. Then
\begin{gather*}\tau(B_0)\le\tau(B_1)\le\cdots\le\tau(B_l).
\end{gather*}
\end{enumerate}
\end{Lemma}

\begin{proof}
(i) If $\tau(B')$ is the minimal triple in the $\Gamma$-orbit of $\tau(B)$ it satisf\/ies~(M1), so by Theo\-rem~\ref{cyclic}(iii) the root is the unique triple satisfying~(M1) and $B'$ must be the root.
For~(ii): $a=b=c=2$ follows easily from Theorem~\ref{cyclic}(ii), and the uniqueness claim follows from the def\/inition of~$\mu_k$.
 (iii)~is obvious if  $\tau(B)=(2,2,2)$. If not, assume~(iii) is false, then there are $1\le j<j'\le l$ such that
 \begin{gather*}
 \tau(B_j)<\tau(B_{j+1})=\tau(B_{j+2})=\cdots=\tau(B_{j'})>\tau(B_{j'+1}).
 \end{gather*}
By (ii) it must be that $j+1=j'$ so that  $\tau(B_j)< \tau(B_{j+1}) > \tau(B_{j+2})$, but $\tau(B_{j+1})$ does not satisfy either~(M1) or~(M2). This is impossible by Theorem~\ref{cyclic}(iii).
\end{proof}

\subsection[Grading on~$\mathcal{A}$]{Grading on $\boldsymbol{\mathcal{A}}$}\label{section4.2}

We now adapt the grading introduced in \cite[Def\/inition 3.1]{G} to the geometric type.
\begin{Definition}
A graded seed is a quadruple $({\bf x},{\bf y},B,G)$ such that
\begin{enumerate}\itemsep=0pt
\item[(i)] $({\bf x},{\bf y},B)$ is a seed of rank $n$, and
\item[(ii)] $G=[g_1,\dots,g_n]^T\in \mathbb{Z}^n$ is an integer column vector such that $BG=0$.
\end{enumerate}
\end{Definition}

Set $\deg_G(x_i)=g_i$ and $\deg\big(x_i^{-1}\big)=-g_i$ for $i\le n$, and set $\deg(y_j)=0$ for all $j$. Extend the grading additively to Laurent monomials hence to the cluster algebra $\mathcal{A}({\bf x},{\bf y},B)$. In~\cite{G} it is proved that under this grading every exchange relation is homogeneous and thus the grading is compatible with mutation.

\begin{Theorem}[\protect{\cite[Corollary 3.4]{G}}]
The cluster algebra $\mathcal{A}({\bf x},{\bf y},B)$ under the above grading is a~$\mathbb{Z}$-graded algebra.
\end{Theorem}

The following two propositions come from the work in \cite{S1} to show that rank three non-acyclic cluster algebras have no maximal green sequences.
\begin{Theorem}[\protect{\cite[Proposition 2.2]{S1}}]
Suppose that $B$ is a $3\times 3$ skew-symmetric non-acyclic matrix. Then the $($column$)$ vector $G=\tau(B)^T$ satisfies $BG=0.$
\end{Theorem}

\begin{Lemma}\label{lemma: grading}
For any graded seed $({\bf x}',B',G')$ in the mutation class of our initial graded seed, we have $G'= \tau(B')^T$.
\end{Lemma}

\begin{proof}
This result follows from \cite[Lemma~2.3]{S1}, but its proof is short, so we reproduce it here. We use induction on mutations. Suppose that $G'= \tau(B')^T$ for a given graded seed $({\bf x}',B',G')$. Let $({\bf x}'',B'',G'')$ be the graded seed obtained by taking mutation $\mu_1$ from $({\bf x}',B',G')$.  Then $x_1''=((x_2')^{|b'_{12}|}+(x_3')^{|b'_{13}|})/x_1'$, so its degree is $g_2' |b'_{12}|-g_1'$. By induction we have
\begin{gather*}
g_2' |b'_{12}|-g_1' =|b'_{31}| |b'_{12}|-|b'_{23}| = |b'_{31}b'_{12}-b'_{23}| =|b''_{23}|,
\end{gather*}
so we get $G''= \tau(B'')^T$. The cases of~$\mu_2$ and~$\mu_3$ are similar.
\end{proof}


In the rest of Section~\ref{section4} we assume that $B$ is a $3\times 3$ skew-symmetric non-acyclic matrix and $G=\tau(B)^T$. In other words, if $\tau(B)=(b,c,a)$ then $\deg(x_1)=b$, $\deg(x_2)=c$, $\deg(x_3)=a$ for any seed $({\bf x}, {\bf y}, B)$. If we look at the quiver associated to our exchange matrix we see that the degree of the cluster variable $x_i$ is the number of arrows between the other two mutable vertices in~$Q_B$. Furthermore, this is a~canonical grading for non-acyclic rank three cluster algebras since regardless of your choice of initial seed the grading imposed on~$\mathcal{A}$ is the same.

\subsection[Construction of an element in $\mathcal{U}\setminus\mathcal{A}$]{Construction of an element in $\boldsymbol{\mathcal{U}\setminus\mathcal{A}}$}\label{section4.3}
Here we shall prove the following main theorem of the paper.

We give $\mathcal{A}$ the $\mathbb{Z}$-grading as in Section~\ref{section4.2}. By Theorem~\ref{cyclic}, we can assume that the initial seed
\begin{gather*}
\Sigma=(\{x_1,x_2,x_3\},{\bf y},B,G),\qquad\textrm{where} \qquad B= \left[ \begin{matrix}
0 & a &-c \\
-a & 0 & b \\
c& -b & 0
\end{matrix}  \right], \qquad G=\left[ \begin{matrix}
b \\ c \\ a  \end{matrix} \right]
\end{gather*}
satisf\/ies $a,b,c\ge2$ and (M1). Then $\deg x_{1}=b$, $\deg x_{2}=c$, $\deg x_{3}=a$, and we let $x_{4},\dots,x_{m}$ be of degree~0.  Furthermore, by permuting the indices if necessary, we assume
$a\ge b\ge c$.


Def\/ine six degree-0 elements as follows
\begin{gather*}
\alpha_i^\pm := \prod_{j=4}^m x_j^{[\pm b_{ji}]_+}, \qquad i=1,2,3.
\end{gather*}
Looking at the seeds neighboring $\Sigma$ we obtain three new cluster variables. Namely,
\begin{gather*}
z_1=\frac{\alpha_1^-x_2^a+\alpha_1^+x_3^c}{x_1},
\qquad z_2=\frac{\alpha_2^+x_1^a+\alpha_2^-x_3^b}{x_2}, \qquad z_3=\frac{\alpha_3^-x_1^c+\alpha_3^+x_2^b}{x_3},
 \end{gather*} which have degree $ac-b$, $ab-c$, and $bc-a$, respectively.

We also use the following theorem, whose proof in~\cite{MM} is for coef\/f\/icient free cluster algebras but clearly applies to cluster algebras of geometric type.
\begin{Theorem}[\protect{\cite[Proposition 6.1.2]{MM}}]
\label{coprime}
If $\mathcal{A}$ is a rank three skew-symmetric non-acyclic cluster algebra, then~$\mathcal{A}$ is totally coprime.
\end{Theorem}

Now we consider a special element in $\mathbb{Q}(x_1,x_2,x_3)$:
\begin{gather*}
Y:=\tilde{x}[(1,0,1)]/x_2^b =\frac{\alpha_1^-\alpha_3^-x_1^cx_2^{a-b}+\alpha_1^-\alpha_3^+x_2^a+\alpha_1^+\alpha_3^+x_3^c}{x_1x_3}.
\end{gather*}
We shall prove Theorem~\ref{main-cor} by showing that the element constructed above is in $\mathcal{U}\setminus\mathcal{A}$.

\begin{proof}[Proof of Theorem \ref{main-cor}]
We f\/irst show that $Y\in \mathcal{U}$.
Note that $\mathcal{A}$ is a rank 3 cluster algebra so it is totally coprime by Theorem~\ref{coprime}. It then suf\/f\/ices to show that $Y\in \mathcal{U}_{\bf x}$ by Theorem~\ref{UB}. Clearly $Y\in\mathbb{ZP}[x_1^{\pm 1},x_2^{\pm1},x_3^{\pm 1}]$, where
$\mathbb{ZP}=\mathbb{Z}[x_4^{\pm1},\dots,x_m^{\pm1}]$. Also,
\begin{gather*}
Y=\frac{\alpha_3^+z_1^c+\alpha_1^-\alpha_3^-\big(\alpha_1^-x_2^a+\alpha_1^+x_3^c\big)^{c-1} x_2^{a-b}}{z_1^{c-1}x_3} \in \mathbb{ZP}\big[z_1^{\pm 1},x_2^{\pm 1},x_3^{\pm 1}\big],\\ Y=\frac{\alpha_1^-\alpha_3^-x_1^c\!\big(\alpha_2^+x_1^a\!+\!\alpha_2^-x_3^b\big)^{a{-}b}\!z_2^b\!+\!\alpha_1^+\alpha_3^+z_2^ax_3^c\!+\!\alpha_1^-\alpha_3^+
\!\big(\alpha_2^+x_1^a\!+\!\alpha_2^-x_3^b\big)^{a}}{x_1z_2^ax_3}
%\\ \hphantom{Y=}{}  
 \in  \mathbb{ZP}\big[x_1^{\pm 1}\!,z_2^{\pm 1}\!,x_3^{\pm 1}\big],\\
  Y=\frac{\alpha_1^-x_2^{a-b}z_3^c+\alpha_1^+\alpha_3^+\big(\alpha_3^-x_1^c+\alpha_3^+x_2^b\big)^{c-1}}{x_1z_3^{c-1}} \in\mathbb{ZP}\big[x_1^{\pm 1},x_2^{\pm 1},z_3^{\pm 1}\big].
\end{gather*}
Therefore we conclude that $Y\in \mathcal{U}_{\bf x}$.



Next, we show that $Y\notin \mathcal{A}$. With respect to our grading of $\mathcal{A}$, $Y$ is homogeneous of degree $ac-b-a.$

Combining Lemmas \ref{lemma:chain inequality} and~\ref{lemma: grading} we have already shown that the degree of cluster variables is non-decreasing as we mutate away from our initial seed. We use this fact to explicitly prove that all cluster variables, except possibly $x_1$, $x_2$, $x_3$, $z_3$, have degree strictly larger than $ac-b-a$.  Indeed, let $x$ be a cluster variable such that $x\neq x_1,x_2,x_3,z_3$. Then $x$ can be written as
\begin{gather*}
x=\mu_{i_l}\cdots\mu_{i_2}\mu_{i_1}(x_k),\quad i_1,\dots,i_l,k\in\{1,2,3\}, \qquad \text{and}\qquad i_s\neq i_{s+1} \quad \text{for} \ \ 1\le s\le l-1.
\end{gather*}
Let $B_0=B$, $B_j=\mu_{i_j}(B_{j-1})$ for $j=1,\dots,l$. Let $\tau(B_j)=(b_j,c_j,a_j)$ for $j=0,\dots,l$. By Lemma~\ref{lemma:chain inequality}, we have that
\begin{gather*}
(b,c,a)=(b_0,c_0,a_0)\le(b_1,c_1,a_1)\le\cdots\le(b_l,c_l,a_l).
\end{gather*}
We prove the claim in the following f\/ive cases.

\textit{Case $k=1$}. We let $r>0$ be the smallest integer such that $i_r=k$ (which exists since $x\neq x_1,x_2,x_3$). It suf\/f\/ices to prove that the degree of $w=\mu_{i_r}\cdots\mu_{i_2}\mu_{i_1}(x_k)$ is larger than $ac-b-a$. Indeed,
\begin{gather*}
\deg w=b_r=a_{r-1}c_{r-1}-b_{r-1}=a_{r-1}c_{r-1}-b\ge ac-b>ac-b-a.
\end{gather*}

\textit{Case $k=2$}. The above proof still works:
\begin{gather*}
\deg w=c_r=a_{r-1}b_{r-1}-c_{r-1}=a_{r-1}b_{r-1}-c\ge ab-c>ac-b-a.
\end{gather*}

\textit{Case  $k=3$, $i_1=1$}. We let $r>0$ be the smallest integer such that $i_r=k$. Then $(b_1,c_1,a_1)=(ac-b,c,a)$, and
\begin{gather*}
\deg w=a_r=b_{r-1}c_{r-1}-a_{r-1}=b_{r-1}c_{r-1}-a\ge b_1c_1-a=(ac-b)c-a>ac-b-a.
\end{gather*}

\textit{Case  $k=3$, $i_1=2$}. Similar to the above case, $(b_1,c_1,a_1)=(b,ab-c,a)$,
\begin{gather*}
\deg w=a_r\ge b_1c_1-a=b(ab-c)-a\ge b^2(a-1)-a>c(a-1)-a\ge ac-b-a.
\end{gather*}
 The second and fourth inequality follow from our assumption that $b \geq c$.

\textit{Case  $k=3$, $i_1=3$}. Then $(b_1,c_1,a_1)=(b,c,bc-a)$, $i_2\neq 3$. We let $r\ge3$ be the smallest integer such that $i_r=k$ (which exists since $x\neq z_3$).  It suf\/f\/ices to prove that the degree of $w=\mu_{i_r}\cdots\mu_{i_2}\mu_{i_1}(x_k)$ is larger than $ac-b-a$.

If $i_2=1$, then $(b_2,c_2,a_2)=(c(bc-a)-b,c,bc-a)$,
\begin{gather*}
\begin{split}
&\deg w =a_r=b_{r-1}c_{r-1}-(bc-a)\ge b_2c_2-(bc-a)=(c(bc-a)-b)c-(bc-a)
\\
& \hphantom{\deg w }{} =(c^2-2)(bc-a)-a\ge(c^2-2)a-a=a(c^2-3)\ge a(c-1)>ac-b-a.
\end{split}
\end{gather*}

If $i_2=2$, then $(b_2,c_2,a_2)=(b,(bc-a)b-c,bc-a)$,
\begin{gather*}
\deg w\ge b((bc-a)b-c)-(bc-a)\ge (c(bc-a)-b)c-(bc-a)>ac-b-a.
\end{gather*}
The second inequality above follows from our assumption that $b \geq c$. One application gives the inequality $(bc-a)b-c\ge c(bc-a)-b$ and a second application replaces a factor of $b$ in the f\/irst term with a factor of~$c$. The last inequality follows from our computation in the previous case.

This completes the proof of the claim that all cluster variables, except possibly $x_1$, $x_2$,~$x_3$,~$z_3$, have degree strictly larger than $ac-b-a$.

Therefore it is suf\/f\/icient to check that $Y$ cannot be written as a linear combination of products of the cluster variables $x_1$, $x_2$, $x_3$, and~$z_3$ since all other cluster variable have larger degree than~$Y$. This is clear as the numerator of~$Y$ is irreducible and none of these cluster variables have a~factor of~$x_1$ in the denominator.
\end{proof}
